% \begin{abstract}
% Time series (TS) reasoning models (TSRMs) have shown promising capabilities in general domains such as weather and traffic, yet they consistently fail on financial TS, which exhibit unique market characteristics. 
% We argue that building a \textbf{Financial TSRM} requires defining core capabilities that span four categories: we construct a $2 \times 2$ taxonomy crossing 1) \textit{single-stock vs.\ multi-stock} analysis with 2) \textit{assessment} of the current state and \textit{prediction} of future behavior, and instantiate them as ten financial reasoning tasks, forming a benchmark called \textbf{FinTSR-Bench} constructed from 250 S\&P~500 stocks (2010--2025).
% As assessment is deterministic (i.e., computable from observable data) while prediction is inherently stochastic (i.e., subject to unobservable external factors), we design distinct CoT strategies for each.
% For assessment tasks, we employ \textbf{Compute-in-CoT}, a programmatic chain-of-thought (CoT) that teaches models to compute answers directly from raw prices.
% For prediction tasks, we propose \textbf{Scenario-Aware CoT}, which generates base-case, adverse, and favorable scenarios before making a probabilistic judgment, mirroring how financial analysts reason under uncertainty.
% The resulting model, \textbf{FinSTaR} (\textbf{Fin}ancial \textbf{S}eries \textbf{T}hinking \textbf{a}nd \textbf{R}easoning), achieves 78.9\% average accuracy on our benchmark, substantially outperforming zero-shot LLM baselines, TSRMs, and TS forecasting models.
% Through systematic ablation, we demonstrate that assessment capabilities serve as the foundation for prediction, and that scenario-aware reasoning provides consistent gains on prediction tasks over standard CoT.
% \end{abstract}
\vspace{-8pt}
\begin{abstract}
\vspace{-8pt}
Time series (TS) reasoning models (TSRMs) have shown promising capabilities in general domains, yet they consistently fail 
%on financial TS, 
on financial domain, 
which exhibit unique characteristics. 
We propose a general $2 \times 2$ capability taxonomy for TSRMs by crossing 1) \textit{single-entity vs.\ multi-entity} analysis with 2) \textit{assessment} of the current state vs.\ \textit{prediction} of future behavior.
We instantiate this taxonomy in the financial domain---where the distinction between deterministic assessment and stochastic prediction is particularly critical---as ten financial reasoning tasks, forming the \textbf{FinTSR-Bench} benchmark based on S\&P stocks.
%(2010--2025).
To this end, we propose \textbf{FinSTaR} (\textbf{Fin}ancial Time \textbf{S}eries \textbf{T}hinking \textbf{a}nd \textbf{R}easoning), 
trained on FinTSR-Bench with distinct chain-of-thought (CoT) strategies tailored to each category.
For \textit{assessment}, which is \textit{deterministic} (i.e., computable from observable data), we employ \textit{Compute-in-CoT}, a programmatic CoT that enables models to derive answers directly from raw prices. 
For \textit{prediction}, which is inherently \textit{stochastic} (i.e., subject to unobservable factors), we adopt \textit{Scenario-Aware CoT}, which generates diverse scenarios before making a judgment, mirroring how financial analysts reason under uncertainty.
The proposed method achieves 78.9\% average accuracy 
on FinTSR-Bench,
%on our benchmark, 
substantially outperforming 
%zero-shot LLM baselines, TSRMs, and TS forecasting models.
LLM and TSRM baselines.
Furthermore, we show that 
the 
four capability categories are complementary and mutually reinforcing through joint training, and that Scenario-Aware CoT consistently improves prediction accuracy over standard CoT.
Code is publicly available at: https://github.com/seunghan96/FinSTaR.
\end{abstract}
\vspace{-10pt}
